% Generated draft: requires mathematical and source review.
\documentclass[11pt]{article}
\usepackage[margin=25mm]{geometry}
\usepackage{fontspec}
\setmainfont{Latin Modern Roman}
\usepackage{amsmath,amssymb,mathtools,unicode-math}
\setmathfont{Latin Modern Math}
\usepackage{enumitem,needspace,xurl,hyperref}
\setlength{\parindent}{0pt}
\setlength{\parskip}{5pt}
\setlength{\emergencystretch}{3em}
\newcommand{\N}{\mathbb{N}}
\newcommand{\Z}{\mathbb{Z}}
\newcommand{\Q}{\mathbb{Q}}
\newcommand{\R}{\mathbb{R}}
\newcommand{\C}{\mathbb{C}}
\newcommand{\F}{\mathbb{F}}
\newcommand{\A}{\mathbb{A}}
\newcommand{\PP}{\mathbb{P}}
\newcommand{\E}{\mathbb{E}}
\newcommand{\eps}{\varepsilon}
\newcommand{\dd}{\,\mathrm{d}}
\begin{document}
\section*{UnsolvedMath ID: 6200105}
\textbf{Source problem code:} catalogue entry
\subsection{Definitions and mathematical statement}
For a finitely generated group \(G\), a subgroup \(H\leq G\), and a finite symmetric generating set \(S\) of \(G\), let \(\Gamma_{G/H,S}\) have vertices the left cosets \(gH\) and edges from \(gH\) to \(sgH\), for \(s\in S\). Give its edges length one and set
\[
 b_{G/H,S}(r)=\#\{gH:d_{\Gamma}(H,gH)\leq r\}\qquad(r\in\N).
\]
Here the source's term \emph{cogrowth} means this Schreier graph volume growth. It does not mean the number of words returning to \(H\). Changing the finite generating set preserves its coarse growth type. Bounded growth means \(\sup_r b(r)<\infty\); linear growth means \(cr\leq b(r)\leq Cr\) for some \(c,C>0\) and all sufficiently large \(r\); exponential growth means \(b(r)\geq ca^r\) for some \(c>0\), \(a>1\), and all sufficiently large \(r\). A finite degree bound supplies an exponential upper bound.

The source asks to compute this growth for notable inclusions \(H\leq G\), in particular:
\begin{enumerate}
\item Prove exponential growth for \(\operatorname{SL}(n,\Z)\leq\operatorname{SL}(n+1,\Z)\), using the usual block inclusion \(A\mapsto\operatorname{diag}(A,1)\). The intended range of \(n\) is not specified.
\item Compute the growth for special subgroups of Coxeter groups. With the customary convention, a special subgroup of a Coxeter system \((W,S)\) is \(W_T=\langle T\rangle\), \(T\subseteq S\); the graph in question is \(\Gamma_{W/W_T,S}\). Finite \(S\) is required for the finite degree volume growth convention above. The source does not specify whether an exact growth function or only its growth type is requested.
\item If \(\Gamma_{G/H,S}\) is Gromov-hyperbolic, must \(b_{G/H,S}\) have bounded, linear, or exponential growth? Gromov-hyperbolic means that geodesic triangles are uniformly thin in the graph metric. The proposed trichotomy excludes all other unbounded subexponential growth types, including higher degree polynomial growth.
\end{enumerate}
The finite generation and growth conventions above make the usual interpretation explicit; they are not stated in the extracted source and require source review.
\subsection{Short English statement}
Determine the volume growth of subgroup Schreier graphs, particularly for the standard inclusion of consecutive integral special linear groups and for special Coxeter subgroups. Decide whether hyperbolicity of such a graph forces bounded, linear, or exponential growth.
\subsection{Sources}
\begin{itemize}
\item Local UnsolvedMath registry, permanent ID 6200105, source code catalogue entry, dated update 2026-08-06; proposed year 2005; no individual proposer supplied.
\item Kapovich, \emph{Problems on Boundaries of Groups and Kleinian Groups} (2005), Problem 105, PDF page 26: \url{https://www.math.ucdavis.edu/~kapovich/EPR/problems.pdf}. The registry reports access on 2026-07-29 and normalized PDF text extraction; the source was not independently checked for this draft.
\item The supplied summary mentions Grigorchuk's cogrowth theory and amenability connections, but gives no precise additional reference. It labels the problem partially solved; the background separately marks present status and statement reuse terms as needing review for this public author-hosted PDF.
\end{itemize}
\end{document}
